\section{Benchmarking}
\label{sec:benchmark}

\begin{figure}[t]
  \centering
  \begin{subfigure}[t]{0.48\linewidth}
    \includegraphics[width=\linewidth]{figures/scaling_overall_n_boxes_grip1x.pdf}
    \caption{$k=1$}
  \end{subfigure}\hfill
  \begin{subfigure}[t]{0.48\linewidth}
    \includegraphics[width=\linewidth]{figures/scaling_overall_n_boxes_grip2x.pdf}
    \caption{$k=2$}
  \end{subfigure}

  \medskip

  \begin{subfigure}[t]{0.48\linewidth}
    \includegraphics[width=\linewidth]{figures/scaling_overall_n_boxes_grip3x.pdf}
    \caption{$k=3$}
  \end{subfigure}
  \caption{Solve time vs. number of boxes, split by gripper multiplier, with
  $S = k \cdot \texttt{boxes\_x\_dim}$. The $y$-axis is logarithmic.}
  \label{fig:scaling-per-grip}
\end{figure}

The benchmark evaluates the MiniZinc model on every
\((\text{pallet}, \text{pack})\) combination obtained from the
pallet presets in Table~\ref{tab:bench-pallets} and the pack sizes in
Table~\ref{tab:bench-packs}.

\begin{table}[h]
  \centering

  \begin{minipage}[t]{0.49\linewidth}
    \centering
    \begin{tabularx}{\linewidth}{lX}
      \toprule
      \textbf{Pallet preset} & \textbf{Dimensions (mm)} \\
      \midrule
      Europallet      & $1200 \times 800$  \\
      Industrie       & $1200 \times 1000$ \\
      US 40$\times$48 & $1219 \times 1016$ \\
      Half-pallet     & $800  \times 600$  \\
      Australian      & $1165 \times 1165$ \\
      \addlinespace
      \bottomrule
    \end{tabularx}
    \caption{Pallet presets used in the benchmark.}
    \label{tab:bench-pallets}
  \end{minipage}
  \hfill
  \begin{minipage}[t]{0.49\linewidth}
    \centering
    \begin{tabularx}{0.8\linewidth}{X}
      \toprule\arraybackslash\textbf{Box size (mm)} \\
      \midrule
      \arraybackslash $50 \times 50$   \\
      \arraybackslash $70 \times 70$   \\
      \arraybackslash $100 \times 100$ \\
      \arraybackslash $120 \times 80$  \\
      \arraybackslash $150 \times 150$ \\
      \arraybackslash $200 \times 200$ \\
      \bottomrule
    \end{tabularx}
    \caption{Box sizes used in the benchmark.}
    \label{tab:bench-packs}
  \end{minipage}
\end{table}
For each combination, the model is tested with four solvers:
\texttt{gecode}, \texttt{cp-sat}, \texttt{chuffed}, and \texttt{highs}.
These solvers represent and solve the problem in different ways.
\texttt{gecode} and \texttt{chuffed} mainly rely on constraint propagation,
\texttt{cp-sat} uses a SAT-based representation, and
\texttt{highs} uses a linear programming representation. Each run is executed
with a timeout of 300 seconds.
{\color{red} SISTEMA DA QUI}
The gripper size is varied as an integer multiple of the box dimension along the
\(x\)-axis: $S \in \{1,2,3\} \cdot \texttt{boxes\_x\_dim}.$
This variation is used to study how the gripper size affects the feasibility and
runtime of the scheduling model.

In the model, this makes \texttt{max\_pkble\_boxes} equal to \(k+1\).
The three benchmark cases therefore correspond to maximum group lengths of
two, three, and four boxes along a pallet row.

All timing plots in this section use a logarithmic scale for the time axis.
This is useful because the benchmark spans several orders of magnitude: some
\texttt{gecode} runs are solved in milliseconds, while several \texttt{highs}
runs reach the 300-second timeout. With a linear scale, the fast solvers would
be compressed near zero and the comparison would be much harder to read.

\subsection{Scaling by gripper size}

The first comparison keeps the gripper multiplier fixed and studies how the
reported solve time changes with the total number of boxes. This separates the
three different grouping regimes before aggregating them into a single plot.
Figure~\ref{fig:scaling-per-grip} shows the resulting scaling curves for
$k=1$, $k=2$, and $k=3$.

The split shows the expected effect of the gripper size. Larger grippers reduce
the number of groups that must be scheduled, and the average number of groups
drops from about $75$ for $k=1$ to about $40$ for $k=3$. As a result, the
curves generally move downward when the multiplier increases.

\texttt{gecode} stays in the lowest band and is almost flat across the whole
range. \texttt{cp-sat} scales more visibly with the number of boxes, but still
remains well below the timeout. \texttt{chuffed} is usually fast, with one
timeout on the densest $k=1$ instance. \texttt{highs} is consistently the
slowest solver.

\subsection{Overall solver behaviour}

After separating the grip sizes, the full benchmark can be viewed as one
aggregate scaling plot. Figure~\ref{fig:scaling-overall} shows solve time with
respect to the total number of boxes on the pallet for every solver across all
instances.

\begin{figure}[t]
  \centering
  \includegraphics[width=0.7\linewidth]{figures/scaling_overall_n_boxes.pdf}
  \caption{Solve time vs. total number of boxes --- all solvers, all
  instances. The $y$-axis is logarithmic.}
  \label{fig:scaling-overall}
\end{figure}

\subsection{Per-pallet breakdown}

Splitting the data by pallet preset, as shown in
Figure~\ref{fig:scaling-per-pallet}, isolates the effect of pack count from
pallet shape. For a fixed pallet, the curves are functions of pack count and
gripper size; differences between pallets reflect the impact of
\texttt{boxes\_along\_x} and \texttt{boxes\_along\_y} on the number of groups.

The larger pallets expose the same behavior seen in the aggregate plots more
clearly: dense layouts with many small boxes are the hard cases, while larger
boxes produce fewer groups and remain easier for all solvers. The log scale is
especially useful in these plots because it shows the sub-second propagation
runs and the near-timeout \texttt{highs} runs in the same figure.

\begin{figure}[t]
  \centering
  \begin{subfigure}[t]{0.45\linewidth}
    \includegraphics[width=\linewidth]{figures/scaling_europallet_n_boxes.pdf}
    \caption{Europallet}
  \end{subfigure}\hfill
  \begin{subfigure}[t]{0.45\linewidth}
    \includegraphics[width=\linewidth]{figures/scaling_industrie_n_boxes.pdf}
    \caption{Industrie}
  \end{subfigure}

  \medskip

  \begin{subfigure}[t]{0.45\linewidth}
    \includegraphics[width=\linewidth]{figures/scaling_us_40x48_n_boxes.pdf}
    \caption{US 40$\times$48}
  \end{subfigure}\hfill
  \begin{subfigure}[t]{0.45\linewidth}
    \includegraphics[width=\linewidth]{figures/scaling_halfpallet_n_boxes.pdf}
    \caption{Half-pallet}
  \end{subfigure}

  \medskip

  \begin{subfigure}[t]{0.45\linewidth}
    \includegraphics[width=\linewidth]{figures/scaling_australian_n_boxes.pdf}
    \caption{Australian}
  \end{subfigure}
  \caption{Solve time vs. number of boxes, broken down by pallet preset. The
  $y$-axis is logarithmic.}
  \label{fig:scaling-per-pallet}
\end{figure}

\subsection{Wall time vs.\ solve time}

The gap between Python-measured wall time and solver-reported solve time
captures MiniZinc flattening and inter-process overhead.
Figure~\ref{fig:wall-vs-solve} plots the two against each other; points
above the diagonal indicate runs in which the harness cost dominated the search
itself. In this plot both axes are logarithmic, so the small overheads of the
fast runs and the long timeout-scale runs can be compared in the same view.

This distinction is especially visible for the fastest solver runs: their
reported solve time is very small, so model construction, flattening, and
process startup become a relevant part of the measured wall time. For the slow
runs near the timeout, instead, wall time and solve time are much closer.

\begin{figure}[t]
  \centering
  \includegraphics[width=0.75\linewidth]{figures/wall_vs_solve.pdf}
  \caption{Wall time vs. solve time per run. The diagonal represents
  zero overhead; both axes are logarithmic.}
  \label{fig:wall-vs-solve}
\end{figure}
